\section*{Agradecimientos}

Durante mucho tiempo, era habitual ver a guitarristas sumergidos en un mar de papeles y móviles distintos para poder tocar las canciones que querían en eucaristías, celebraciones, campamentos, etc. Tanto es así, que incluso canciones bien conocidas por todos y que habitualmente se escuchaban en nuestra pastoral no estaban en el cancionero que utilizábamos y al que teníamos tanto cariño. Precisamente de esa necesidad de renovación, nace este nuevo cancionero para todos.

"La mies era mucha, pero los obreros más" sería la oración con la que definiría el nacimiento de este cancionero. El proyecto fue acogido desde el principio con entusiasmo y aunque el trabajo era inmeso, más lo era la ilusión de los que se prestaban. No puedo más que dar las gracias a todos los que han puesto un poco de su tiempo, de su comunidad, de su ilusión o de su música en este cancionero para todos. Desde los que participaron activamente en su gestación, hasta los que seguían desde cerca el proyecto alentando a todos los que se afanaban en construirlo.

Por último, dar gracias a Dios por haber podido ser partícipes y testigos de tan bella aventura. Gracias Señor por proporcionar cuantas manos hicieron falta, incluso cuando la carga era demasiada para ellas. Gracias Señor por la música de este cancionero, que ha hecho acompañado la fe de tantos de nosotros a través de algo tan sencillo como notas en un cuaderno. Envía Señor tu Espíritu. Envíalo para inspirar a cuantos lean, escuchen y vivan el Evangelio que anuncian estas canciones. Envía Señor tu Espíritu para que esta música sea testimonio de las vidas y experiencias de fe de aquellos que las escribieron. Envía Señor tu Espíritu para que ésta sea instrumento tuyo y nos haga acercarnos cada día más a Cristo.

Para terminar, te invito a ti que lees a cantar, porque como dicen las voces populares:
\begin{quote}
    Cantar es orar dos veces
\end{quote}